\chapter{System Architecture}
\section{Architecture Overview}
The project is developed incrementally. The first milestone is a working API connected to MySQL. Document upload, preprocessing, OCR and review follow. Docker Compose is introduced once individual modules work, services are separated where there is a clear operational benefit, and Kubernetes is a later deployment target, not a prerequisite.

\begin{figure}[H]\centering
\resizebox{\textwidth}{!}{\begin{tikzpicture}[x=1cm,y=1cm]
% layer bands
\fill[mist]  (0,8.5) rectangle (16,10.1);
\fill[mint]  (0,6.1) rectangle (16,8.1);
\fill[sand]  (0,3.0) rectangle (16,5.7);
\fill[lilac] (0,0.6) rectangle (16,2.6);
\fill[blush] (0,-1.4) rectangle (16,0.1);
\node[rotate=90,font=\sffamily\scriptsize\bfseries,text=navy] at (0.3,9.3) {PRESENTATION};
\node[rotate=90,font=\sffamily\scriptsize\bfseries,text=navy] at (0.3,7.1) {API LAYER};
\node[rotate=90,font=\sffamily\scriptsize\bfseries,text=navy] at (0.3,4.35) {DOMAIN SERVICES};
\node[rotate=90,font=\sffamily\scriptsize\bfseries,text=navy] at (0.3,1.6) {DATA};
\node[rotate=90,font=\sffamily\scriptsize\bfseries,text=navy] at (0.3,-0.65) {EXTERNAL};
% presentation
\node[boxT,minimum width=7cm,minimum height=1cm] (ui) at (8.2,9.3) {Flutter Web dashboard\\[-1pt]\scriptsize upload $\cdot$ list $\cdot$ review $\cdot$ approve $\cdot$ metrics};
% api
\foreach \x/\t in {2.3/Auth,5.0/Documents,7.9/Review,10.8/ERP,13.7/Health}{
  \node[box,minimum width=2.3cm] (r\t) at (\x+0.4,7.1) {\textbf{\t}\\[-2pt]\scriptsize router};}
\node[font=\sffamily\scriptsize,text=slate] at (8.2,7.85) {FastAPI $\cdot$ Pydantic validation $\cdot$ JWT / RBAC $\cdot$ versioned routes \texttt{/api/v1}};
% domain
\foreach \x/\t/\s in {2.3/Document/service,4.95/Pre-/processing,7.6/OCR/engine,10.25/Extraction/rules,12.9/Validation/+ review,15.0/ERP/adapter}{}
\node[boxA,text width=2.0cm,minimum height=1.3cm] (sdoc) at (2.4,4.35)  {\textbf{Document}\\service};
\node[boxA,text width=2.0cm,minimum height=1.3cm] (spre) at (4.9,4.35)  {\textbf{Preprocessing}\\OpenCV};
\node[boxA,text width=2.0cm,minimum height=1.3cm] (socr) at (7.4,4.35)  {\textbf{OCR engine}\\Tesseract\\+ CRNN};
\node[boxA,text width=2.0cm,minimum height=1.3cm] (sext) at (9.9,4.35)  {\textbf{Extraction}\\rules + layout};
\node[boxA,text width=2.0cm,minimum height=1.3cm] (sval) at (12.4,4.35) {\textbf{Validation}\\+ review};
\node[boxA,text width=2.0cm,minimum height=1.3cm] (serp) at (14.9,4.35) {\textbf{ERP}\\adapter};
% data
\node[boxV,minimum width=3.6cm,minimum height=1.1cm] (fs) at (3.6,1.6) {\textbf{Protected file storage}\\\scriptsize originals $\cdot$ page images};
\node[boxV,minimum width=2.6cm,minimum height=1.1cm,dashed] (q) at (7.1,1.6) {\textbf{Job queue}\\\scriptsize later stage};
\node[boxV,minimum width=4.0cm,minimum height=1.1cm] (db) at (10.8,1.6) {\textbf{MySQL 8}\\\scriptsize users $\cdot$ documents $\cdot$ fields $\cdot$ audit};
% external
\node[boxC,minimum width=4.2cm,minimum height=0.9cm] (erp) at (12.9,-0.65) {ERP\\[-2pt]\scriptsize Odoo or mock API};
% arrows
\draw[arr] (ui.south) -- ++(0,-0.35) -| (rDocuments.north);
\draw[arr] (rDocuments.south) -- ++(0,-0.4) -| (sdoc.north);
\draw[arr] (rReview.south) -- ++(0,-0.4) -| (sval.north);
\draw[arr] (rERP.south) -- ++(0,-0.4) -| (serp.north);
\draw[arrd] (sdoc.south) -- (sdoc.south |- fs.north);
\draw[arrd] (spre.south) -- (spre.south |- fs.north);
\draw[arrd] (sext.south) -- (sext.south |- db.north);
\draw[arrd] (sval.south) -- (sval.south |- db.north);
\draw[arr]  (serp.south) -- (serp.south |- erp.north);
\end{tikzpicture}
}
\caption{Layered logical architecture. Solid arrows are requests; dashed arrows are data access. The job queue is introduced only in a later stage.}\label{fig:arch}
\end{figure}

\section{Logical Modules}
{\tblstart\small
\begin{longtable}{L{3.0cm}L{6.2cm}L{4.7cm}}
\hdrrow \hdr{Module} & \hdr{Responsibilities} & \hdr{Input $\rightarrow$ output}\\ \endhead
React UI & Sign-in, upload, document list, review workspace, status dashboard. & User actions $\rightarrow$ API calls and views.\\
API layer & Request validation, routing, authentication, consistent errors, API versioning. & HTTP/JSON.\\
Authentication & User identity, roles, token lifecycle, account status. & Credentials $\rightarrow$ principal.\\
Document service & Upload validation, metadata, storage reference, job lifecycle. & File $\rightarrow$ document ID and state.\\
Preprocessing / OCR & PDF rendering, image enhancement, line segmentation, recognition. & Page $\rightarrow$ text, regions, confidence.\\
Extraction & Turn OCR output into typed invoice fields and keep source evidence. & Text + layout $\rightarrow$ candidate fields.\\
Validation / review & Business rules, confidence policy, corrections, approvals. & Candidates $\rightarrow$ approved record.\\
ERP adapter & Map approved records, authenticate, submit, reconcile status. & Record $\rightarrow$ ERP response.\\
MySQL & Users, metadata, fields, workflow states, attempts, audit records. & Transactional data.\\
Protected storage & Original uploads and derived page images outside the web root. & Files.\\
\hline
\end{longtable}}

\section{Why Start with a Modular Monolith?}
Microservices add network calls, inter-service authentication, distributed failure modes and deployment effort. For a student project the sound path is to build clearly separated modules inside one backend, test their interfaces, and extract services only where needed (typically the OCR worker first, because it is slow and resource-hungry). This keeps the service-oriented design while the first milestone stays achievable. Chapter~\ref{ch:deploy} shows the Docker Compose and Kubernetes stages.

\section{Repository Structure}
\begin{lstlisting}[language=bash]
intelligent-document-processing/
  backend/
    app/
      main.py
      api/              # routers (health, auth, documents, review, erp)
      core/             # settings, security, logging
      database/         # engine, sessions
      models/           # SQLAlchemy models
      schemas/          # Pydantic schemas
      services/         # documents, preprocessing, ocr, extraction,
                        # validation, integration
      workers/          # background jobs (later)
    alembic/            # migrations
    tests/
    Dockerfile  requirements.txt
  frontend/react_app/   # React (Node.js, Vite)
  ml/                   # datasets (not committed), training, evaluation
  infrastructure/       # docker, compose, kubernetes
  docs/
  docker-compose.yml  .env.example  README.md
\end{lstlisting}

\chapter{UML Modelling}\label{ch:uml}
This chapter presents the system from several standard UML viewpoints. Each diagram answers a different question: who uses the system (use cases), what it stores (class and ER), how a document changes state, how work flows, and how components talk to each other (sequence).

\section{Use-Case Diagram}
\begin{figure}[H]\centering
\resizebox{\textwidth}{!}{\begin{tikzpicture}[x=1cm,y=1cm]
\draw[rounded corners=6pt,fill=mist,draw=navy,line width=0.8pt] (3.0,-0.4) rectangle (14.0,10.9);
\node[font=\sffamily\bfseries\small,text=navy] at (8.5,10.5) {\projname{} system};
\node[ucase,minimum width=4.4cm] (u1) at (6.6,9.3)  {Upload document};
\node[ucase,minimum width=4.4cm] (u2) at (6.6,8.0)  {Track status and search documents};
\node[ucase,minimum width=4.4cm] (u3) at (6.6,6.7)  {Review and correct fields};
\node[ucase,minimum width=4.4cm] (u4) at (6.6,5.4)  {Approve or reject record};
\node[ucase,minimum width=4.4cm] (u5) at (6.6,4.1)  {Submit approved record to ERP};
\node[ucase,minimum width=4.4cm] (u6) at (6.6,2.8)  {View audit trail and metrics};
\node[ucase,minimum width=4.4cm] (u7) at (6.6,1.5)  {Configure rules and thresholds};
\node[ucase,minimum width=4.4cm] (u8) at (6.6,0.2)  {Manage users and roles};
\node[ucase,minimum width=2.3cm,fill=sand,draw=amber!80!black] (v1) at (12.2,9.3) {Validate file};
\node[ucase,minimum width=2.3cm,fill=sand,draw=amber!80!black] (v2) at (12.2,4.1) {Validate ERP payload};
\draw[arrd] (u1) -- node[lbl,above]{\stereo{include}} (v1);
\draw[arrd] (u5) -- node[lbl,above]{\stereo{include}} (v2);
\actor{op}{1.2}{8.1}{Document\\operator}
\actor{rv}{1.2}{4.8}{Reviewer /\\accountant}
\actor{ad}{1.2}{1.0}{Administrator /\\supervisor}
\node[box,fill=blush,draw=coral,minimum width=1.9cm,minimum height=1.4cm,align=center] (erp) at (15.9,4.1) {\stereo{system}\\\textbf{ERP / Odoo}\\(or mock API)};
\foreach \a/\u in {op/u1,op/u2, rv/u2,rv/u3,rv/u4, ad/u6,ad/u7,ad/u8, rv/u6}{\draw[color=navy,line width=0.6pt] (\a) -- (\u);}
\draw[color=navy,line width=0.6pt] (erp.west) -- (14.0,4.1) -- (v2.east);
\end{tikzpicture}
}
\caption{Use cases and actors. Every use case requires the user to be signed in (FR-01). \stereo{include} marks behaviour that is always executed as part of another use case.}\label{fig:usecase}
\end{figure}

\subsection*{Use-case summary}
{\tblstart\small
\begin{longtable}{L{4.2cm}L{3.1cm}L{6.5cm}}
\hdrrow \hdr{Use case} & \hdr{Actor} & \hdr{Outcome}\\ \endhead
Upload document & Operator & A valid file is stored, a document ID is returned, a job is queued.\\
Track status and search & Operator, Reviewer & Current state and history of each permitted document.\\
Review and correct fields & Reviewer & Corrected values saved with reasons and audit events.\\
Approve or reject record & Reviewer & Record becomes APPROVED or REJECTED.\\
Submit approved record to ERP & ERP system & ERP reference stored; failures recorded and retryable.\\
View audit trail and metrics & Administrator, Reviewer & Traceable history and processing statistics.\\
Configure rules and thresholds & Administrator & Validation rules and review thresholds updated.\\
Manage users and roles & Administrator & Accounts created, deactivated, roles assigned.\\
\hline
\end{longtable}}

\section{Class Diagrams}
\subsection{Domain model}
\begin{figure}[H]\centering
\resizebox{\textwidth}{!}{\begin{tikzpicture}[x=1cm,y=1cm]
\node[uclass] (user) at (0,0) {\cname{User}\nodepart{two}
 +user\_id : UUID\\ +email : str\\ +password\_hash : str\\ +role : Role\\ +is\_active : bool
 \nodepart{three}+authenticate()\\ +has\_permission(p)};
\node[uclass] (run) at (5.8,0) {\cname{ProcessingRun}\nodepart{two}
 +run\_id : UUID\\ +model\_version : str\\ +preproc\_config : JSON\\ +started\_at, ended\_at\\ +error\_category : str
 \nodepart{three}+start()\\ +finish(status)};
\node[uclass] (fld) at (11.6,0) {\cname{ExtractedField}\nodepart{two}
 +field\_id : UUID\\ +name : str\\ +raw\_value, value : str\\ +confidence : float\\ +page, bbox : int[ ]\\ +validation : PASS|REVIEW|FAIL
 \nodepart{three}+correct(v, reason)};
\node[uclass] (doc) at (5.8,-5.0) {\cname{Document}\nodepart{two}
 +document\_id : UUID\\ +storage\_key : str\\ +sha256 : str\\ +mime\_type : str\\ +size\_bytes : int\\ +status : DocumentStatus\\ +created\_at : datetime
 \nodepart{three}+transition(to)\\ +is\_duplicate\_of(d)};
\node[uclass] (rev) at (0,-5.0) {\cname{Review}\nodepart{two}
 +review\_id : UUID\\ +decision : Decision\\ +notes : str\\ +decided\_at : datetime
 \nodepart{three}+approve()\\ +reject(reason)};
\node[uclass] (erp) at (11.6,-5.0) {\cname{ERPAttempt}\nodepart{two}
 +attempt\_id : UUID\\ +idempotency\_key : str\\ +destination : str\\ +status : str\\ +external\_ref : str\\ +error : str
 \nodepart{three}+submit(payload)\\ +retry()};
\node[uclass] (aud) at (5.8,-10.4) {\cname{AuditEvent}\nodepart{two}
 +event\_id : UUID\\ +actor\_id : UUID\\ +document\_id : UUID\\ +action : str\\ +before, after : JSON\\ +at : datetime
 \nodepart{three}+record(action, ctx)};
\node[uenum] (enum) at (11.6,-10.4) {\cname{\stereo{enum} DocumentStatus}\nodepart{two}
 UPLOADED, QUEUED,\\ PROCESSING, NEEDS\_REVIEW,\\ APPROVED, REJECTED,\\ ERP\_SENT, ERP\_FAILED, FAILED};
% relations
\draw[color=navy,line width=0.7pt] (user.south) -- node[lbl,pos=0.15,right]{1} node[lbl,pos=0.5,right]{reviews} node[lbl,pos=0.85,right]{0..*} (user.south |- rev.north);
\draw[color=navy,line width=0.7pt] (user.south east) -- node[lbl,pos=0.12,above right]{1} node[lbl,pos=0.5,sloped,above]{uploads} node[lbl,pos=0.88,above]{0..*} (doc.north west);
\draw[color=navy,line width=0.7pt] (doc.north) -- node[lbl,pos=0.1,left]{1} node[lbl,pos=0.5,left]{processed by} node[lbl,pos=0.9,left]{1..*} (run.south);
\draw[color=navy,line width=0.7pt,-{Diamond[open,length=3mm]}] (fld.west) -- node[lbl,pos=0.12,above]{0..*} node[lbl,pos=0.88,above]{1} (run.east);
\draw[color=navy,line width=0.7pt] (rev.east) -- node[lbl,pos=0.12,above]{0..*} node[lbl,pos=0.88,above]{1} (doc.west);
\draw[color=navy,line width=0.7pt] (doc.east) -- node[lbl,pos=0.12,above]{1} node[lbl,pos=0.88,above]{0..*} (erp.west);
\draw[color=navy,line width=0.7pt] (doc.south) -- node[lbl,pos=0.15,left]{1} node[lbl,pos=0.5,left]{logged in} node[lbl,pos=0.85,left]{0..*} (aud.north);
\draw[color=slate,dashed,line width=0.7pt,-{Stealth[open]}] (doc.south east) -- node[lbl,sloped,above]{uses} (enum.north west);
\end{tikzpicture}
}
\caption{Domain class diagram. A hollow diamond marks composition-like ownership (a processing run owns its extracted fields); plain lines are associations with multiplicities.}\label{fig:classdomain}
\end{figure}

\subsection{Service layer and extension points}
Two interfaces isolate the parts most likely to change. \texttt{OcrEngine} lets the pipeline use Tesseract (baseline) or the custom CRNN without modification, which is exactly what the comparison experiment needs. \texttt{ErpAdapter} keeps ERP-specific code (mock now, Odoo later) out of the document-processing logic (research question~5).
\begin{figure}[H]\centering
\resizebox{\textwidth}{!}{\begin{tikzpicture}[x=1cm,y=1cm]
\node[uclass2] (pipe) at (3.7,0) {\cname{ProcessingPipeline}\nodepart{two}
 +run(document\_id)\\ -preprocess(pages)\\ -recognize(image) : OcrResult\\ -extract(text) : Fields\\ -validate(fields) : Findings};
\node[uclass2] (integ) at (12.3,0) {\cname{IntegrationService}\nodepart{two}
 +submit(document\_id)\\ -build\_payload(record)\\ -idempotency\_key(record)\\ +retry(attempt\_id)};
\node[uiface] (ocr) at (3.7,-3.9) {\cname{\stereo{interface}\ OcrEngine}\nodepart{two}
 +recognize(image) : OcrResult\\ +name : str\\ +version : str};
\node[uiface] (erpa) at (12.3,-3.9) {\cname{\stereo{interface}\ ErpAdapter}\nodepart{two}
 +submit(payload) : ErpResult\\ +get\_status(ref) : Status};
\node[uclass2,text width=3.0cm] (tess) at (1.7,-7.4) {\cname{TesseractEngine}\nodepart{two}
 baseline OCR\\ (comparison)};
\node[uclass2,text width=3.0cm] (crnn) at (5.7,-7.4) {\cname{CrnnEngine}\nodepart{two}
 PyTorch CRNN + CTC\\ model\_version};
\node[uclass2,text width=3.0cm] (mock) at (10.5,-7.4) {\cname{MockErpAdapter}\nodepart{two}
 local test endpoint\\ simulates success/failure};
\node[uclass2,text width=3.0cm] (odoo) at (14.4,-7.4) {\cname{OdooAdapter}\nodepart{two}
 maps fields to Odoo\\ vendor bill (verified later)};
\draw[arrd] (pipe.south) -- node[lbl,right]{uses} (ocr.north);
\draw[arrd] (integ.south) -- node[lbl,right]{uses} (erpa.north);
\foreach \c/\i in {tess/ocr,crnn/ocr,mock/erpa,odoo/erpa}{
  \draw[color=slate,dashed,line width=0.7pt,-{Triangle[open,length=3.2mm,width=3mm]}] (\c.north) -- ++(0,0.6) -| (\i.south);}
\node[font=\sffamily\scriptsize,text=slate,text width=6.2cm,align=center] at (8.0,-9.4) {Dashed arrows with open heads show \emph{realization}: concrete engines and adapters implement the same interface, so either can be swapped without changing the pipeline.};
\end{tikzpicture}
}
\caption{Strategy pattern for OCR engines and ERP adapters.}\label{fig:classservices}
\end{figure}

\section{State Machine of a Document}
A document moves through explicit states. The implementation must reject invalid transitions (for example, an unreviewed document must never jump to \texttt{ERP\_SENT}).
\begin{figure}[H]\centering
\resizebox{0.95\textwidth}{!}{\begin{tikzpicture}[x=1cm,y=1cm]
\node[initial] (i) at (-0.6,0) {};
\node[state] (up) at (1.5,0) {UPLOADED};
\node[state] (qu) at (6.0,0) {QUEUED};
\node[state] (pr) at (10.5,0) {PROCESSING};
\node[stateC] (fa) at (6.0,-3.2) {FAILED};
\node[stateA] (nr) at (10.5,-3.2) {NEEDS\_REVIEW};
\node[stateC] (rj) at (14.6,-3.2) {REJECTED};
\node[stateV] (ap) at (10.5,-6.4) {APPROVED};
\node[state] (es) at (8.0,-9.6) {ERP\_SENT};
\node[stateC] (ef) at (13.0,-9.6) {ERP\_FAILED};
\node[final] (f) at (8.0,-11.6) {};
\node[final] (f2) at (14.6,-5.2) {};
\draw[arr] (i) -- (up);
\draw[arr] (up) -- node[lbl,above]{create job} (qu);
\draw[arr] (qu) -- node[lbl,above]{picked up} (pr);
\draw[arr] (pr) -- node[lbl,right]{OCR, extraction, validation} (nr);
\draw[arr] (pr.south west) -- node[lbl,above,sloped,pos=0.45]{error} (fa.north east);
\draw[arr] (fa) -- node[lbl,left]{retry} (qu);
\draw[arr] (nr) -- node[lbl,above]{reject} (rj);
\draw[arr] (rj) -- (f2);
\draw[arr] (nr) -- node[lbl,right]{approve} (ap);
\draw[arr] (ap.south) -- ++(0,-0.5) -| node[lbl,pos=0.25,above]{submit OK} (es.north);
\draw[arr] (ap.south) -- ++(0,-0.5) -| node[lbl,pos=0.25,above]{submit fails} (ef.north);
\draw[arr] (ef) -- node[lbl,above]{retry} (es);
\draw[arr] (es) -- (f);
\end{tikzpicture}
}
\caption{Document lifecycle. Initial state is the filled dot; double circles are final states.}\label{fig:state}
\end{figure}

{\tblstart\small
\begin{longtable}{L{3.4cm}L{3.4cm}L{7.0cm}}
\hdrrow \hdr{From} & \hdr{To} & \hdr{Trigger / guard}\\ \endhead
UPLOADED & QUEUED & Upload accepted and stored.\\
QUEUED & PROCESSING & Worker starts a processing run.\\
PROCESSING & NEEDS\_REVIEW & Extraction and validation complete (human approval is mandatory in the prototype).\\
PROCESSING & FAILED & Unrecoverable step error; error category stored.\\
FAILED & QUEUED & Manual or bounded automatic retry (FR-21).\\
NEEDS\_REVIEW & APPROVED / REJECTED & Authorised reviewer decision.\\
REJECTED & QUEUED & Reprocessing requested; creates a new run (FR-21).\\
APPROVED & ERP\_SENT / ERP\_FAILED & ERP response.\\
ERP\_FAILED & ERP\_SENT & Retry with the same idempotency key.\\
\hline
\end{longtable}}

\section{Activity Diagram}
\begin{figure}[H]\centering
\resizebox{!}{0.66\textheight}{\begin{tikzpicture}[x=1cm,y=1cm]
% lanes
\foreach \x/\t/\c in {0/Operator/mist, 4/System (API + worker)/mint, 8/Reviewer/sand, 12/ERP adapter/blush}{
  \fill[\c] (\x,-16.4) rectangle (\x+4,0.2);
  \draw[navy,line width=0.5pt] (\x,-16.4) rectangle (\x+4,0.2);
  \fill[navy] (\x,-0.6) rectangle (\x+4,0.2);
  \node[text=white,font=\sffamily\bfseries\scriptsize] at (\x+2,-0.2) {\t};}
% nodes
\node[initial] (s) at (2,-1.2) {};
\node[box,minimum width=2.8cm] (a1) at (2,-2.3) {Select and upload\\ PDF / PNG / JPEG};
\node[box,minimum width=2.8cm] (a2) at (6,-2.3) {Check type, size,\\ integrity, permissions};
\node[diamond,aspect=2,draw=navy,fill=white,inner sep=1pt,font=\sffamily\scriptsize,align=center] (d1) at (6,-4.0) {Accepted?};
\node[box,minimum width=2.8cm,fill=blush,draw=coral] (a3) at (2,-4.0) {Show rejection\\ reason};
\node[final] (e1) at (2,-5.3) {};
\node[box,minimum width=2.8cm,text width=3.1cm] (a4) at (6,-5.5) {Store original, SHA-256, create job};
\node[box,minimum width=2.8cm] (a5) at (6,-6.8) {Render pages and\\ preprocess images};
\node[box,minimum width=2.8cm] (a6) at (6,-8.1) {Recognize text (OCR)};
\node[box,minimum width=2.8cm] (a7) at (6,-9.4) {Extract invoice fields};
\node[box,minimum width=2.8cm] (a8) at (6,-10.7) {Validate and score\\ confidence};
\node[box,minimum width=2.8cm] (b1) at (10,-10.7) {Inspect evidence,\\ correct fields};
\node[diamond,aspect=2,draw=navy,fill=white,inner sep=1pt,font=\sffamily\scriptsize,align=center] (d2) at (10,-12.3) {Approve?};
\node[box,minimum width=2.8cm,fill=blush,draw=coral] (b2) at (10,-14.0) {Reject with reason};
\node[final] (e2) at (10,-15.1) {};
\node[box,minimum width=2.8cm] (c1) at (14,-12.3) {Map and validate\\ ERP payload};
\node[box,minimum width=2.8cm] (c2) at (14,-13.6) {Submit with\\ idempotency key};
\node[box,minimum width=2.8cm] (c3) at (14,-14.9) {Record response\\ and audit event};
\node[final] (e3) at (14,-15.9) {};
% edges
\draw[arr] (s) -- (a1);
\draw[arr] (a1) -- (a2.west);
\draw[arr] (a2) -- (d1);
\draw[arr] (d1.west) -- node[lbl,above]{No} (a3.east);
\draw[arr] (a3) -- (e1);
\draw[arr] (d1) -- node[lbl,right]{Yes} (a4);
\foreach \a/\b in {a4/a5,a5/a6,a6/a7,a7/a8}{\draw[arr] (\a) -- (\b);}
\draw[arr] (a8) -- (b1);
\draw[arr] (b1) -- (d2);
\draw[arr] (d2) -- node[lbl,right]{No} (b2);
\draw[arr] (b2) -- (e2);
\draw[arr] (d2.east) -- node[lbl,above]{Yes} (c1.west);
\draw[arr] (c1) -- (c2);
\draw[arr] (c2) -- (c3);
\draw[arr] (c3) -- (e3);
\end{tikzpicture}
}
\caption{Swimlane activity diagram from upload to controlled ERP submission. Rejection ends the workflow without any ERP call.}\label{fig:activity}
\end{figure}

\section{Sequence Diagrams}
\begin{figure}[H]\centering
\resizebox{\textwidth}{!}{\begin{tikzpicture}[x=1cm,y=1cm]
\def\ys{-0.9}
\foreach \x/\n/\c in {0.9/Operator/mist, 3.7/Flutter UI/mint, 6.5/FastAPI/mint, 9.3/Document\\service/sand, 12.1/Worker\\(OCR)/sand, 14.9/MySQL/lilac}{
  \node[box,fill=\c,minimum width=2.3cm,text width=2.1cm,minimum height=0.9cm] (p\x) at (\x,0) {\textbf{\n}};
  \draw[lifeline] (\x,-0.45) -- (\x,-12.8);
  \fill[navy!15] (\x-0.08,-1.2) rectangle (\x+0.08,-12.2);}
% messages (y descending)
\draw[msg] (0.9,-1.4) -- node[lbl,above]{select file} (3.7,-1.4);
\draw[msg] (3.7,-2.2) -- node[lbl,above]{POST /api/v1/documents (JWT)} (6.5,-2.2);
\draw[msg] (6.5,-3.0) .. controls (7.6,-3.0) and (7.6,-3.5) .. node[lbl,right,pos=0.5]{authenticate, authorize} (6.58,-3.5);
\draw[msg] (6.5,-4.3) -- node[lbl,above]{validate + store} (9.3,-4.3);
\draw[msg] (9.3,-5.1) -- node[lbl,above]{INSERT document, job = QUEUED} (14.9,-5.1);
\draw[ret] (9.3,-5.9) -- node[lbl,above]{document\_id} (6.5,-5.9);
\draw[ret] (6.5,-6.5) -- node[lbl,above]{202 Accepted \{document\_id\}} (3.7,-6.5);
\draw[msg] (9.3,-7.3) -- node[lbl,above]{enqueue / start job} (12.1,-7.3);
\draw[msg] (12.1,-8.1) .. controls (13.2,-8.1) and (13.2,-8.9) .. node[lbl,right,pos=0.5]{preprocess, OCR, extract, validate} (12.18,-8.9);
\draw[msg] (12.1,-9.7) -- node[lbl,above]{save fields, status = NEEDS\_REVIEW} (14.9,-9.7);
\draw[msg] (3.7,-10.5) -- node[lbl,above]{GET /documents/\{id\} (poll)} (6.5,-10.5);
\draw[msg] (6.5,-11.1) .. controls (10.7,-11.1) and (13.2,-11.1) .. node[lbl,above,pos=0.5]{SELECT status, fields} (14.9,-11.1);
\draw[ret] (6.5,-11.7) -- node[lbl,above]{status + fields} (3.7,-11.7);
\draw[ret] (3.7,-12.2) -- node[lbl,above]{show result} (0.9,-12.2);
% async note
\node[draw=amber!80!black,fill=sand,font=\sffamily\scriptsize,align=left,inner sep=3pt] at (7.2,-8.5) {OCR runs asynchronously:\\ the upload request\\ returns before OCR ends.};
\end{tikzpicture}
}
\caption{Upload to OCR result. The upload request returns before OCR finishes; the UI polls the status endpoint.}\label{fig:sequpload}
\end{figure}

\begin{figure}[H]\centering
\resizebox{\textwidth}{!}{\begin{tikzpicture}[x=1cm,y=1cm]
\foreach \x/\n/\c in {0.9/Reviewer/mist, 3.7/Flutter UI/mint, 6.5/FastAPI/mint, 9.3/Review +\\audit/sand, 12.1/ERP\\adapter/sand, 14.9/ERP\\(Odoo or mock)/blush}{
  \node[box,fill=\c,minimum width=2.3cm,text width=2.1cm,minimum height=0.9cm] at (\x,0) {\textbf{\n}};
  \draw[lifeline] (\x,-0.45) -- (\x,-13.6);
  \fill[navy!15] (\x-0.08,-1.2) rectangle (\x+0.08,-13.0);}
\draw[msg] (0.9,-1.4) -- node[lbl,above]{edit field + reason} (3.7,-1.4);
\draw[msg] (3.7,-2.1) -- node[lbl,above]{PATCH /documents/\{id\}/fields} (6.5,-2.1);
\draw[msg] (6.5,-2.8) -- node[lbl,above]{validate role, save correction} (9.3,-2.8);
\draw[msg] (9.3,-3.5) .. controls (11,-3.5) and (11,-4.0) .. node[lbl,right,pos=0.5]{AuditEvent: before / after} (9.38,-4.0);
\draw[msg] (0.9,-4.9) -- node[lbl,above]{click Approve} (3.7,-4.9);
\draw[msg] (3.7,-5.6) -- node[lbl,above]{POST .../approve} (6.5,-5.6);
\draw[msg] (6.5,-6.3) -- node[lbl,above]{status = APPROVED} (9.3,-6.3);
\draw[msg] (3.7,-7.2) -- node[lbl,above]{POST .../submit} (6.5,-7.2);
\draw[msg] (6.5,-7.9) -- node[lbl,above]{submit(document\_id)} (12.1,-7.9);
\draw[msg] (12.1,-8.6) .. controls (13.2,-8.6) and (13.2,-9.1) .. node[lbl,right,pos=0.5]{build payload, key} (12.18,-9.1);
\draw[msg] (12.1,-9.8) -- node[lbl,above]{POST invoice + idempotency key} (14.9,-9.8);
% alt frame
\draw[amber!80!black,line width=0.8pt] (11.2,-10.3) rectangle (15.9,-12.8);
\fill[amber!80!black] (11.2,-10.3) rectangle (11.9,-10.65);
\node[text=white,font=\sffamily\bfseries\tiny] at (11.55,-10.47) {alt};
\draw[ret] (14.9,-11.0) -- node[lbl,above]{ok: external\_ref} (12.1,-11.0);
\node[font=\sffamily\tiny,text=slate,fill=white,inner sep=1pt] at (13.5,-11.42) {[success] status = ERP\_SENT};
\draw[dashed,amber!80!black] (11.2,-11.65) -- (15.9,-11.65);
\draw[ret] (14.9,-12.2) -- node[lbl,above]{error} (12.1,-12.2);
\node[font=\sffamily\tiny,text=slate,fill=white,inner sep=1pt] at (13.5,-12.58) {[failure] ERP\_FAILED, retry later};
\draw[ret] (12.1,-13.0) -- node[lbl,above]{result} (3.7,-13.0);
\end{tikzpicture}
}
\caption{Review, approval and ERP submission, including the success and failure branches of the integration call.}\label{fig:seqreview}
\end{figure}
